\documentclass[a4paper]{article}
\usepackage{times}
\usepackage[utf8]{inputenc}
\usepackage[T1]{fontenc}
\usepackage{selinput}
\usepackage{upquote}
\usepackage[margin=2cm, rmargin=4cm, tmargin=3cm]{geometry}
\usepackage{tcolorbox}
\usepackage{xspace}
\usepackage{amsmath}
\usepackage{float}
\usepackage[french]{babel}
\usepackage{hyperref}
\usepackage{xurl}
\usepackage{fontawesome6}
\usepackage{enumitem}
\usepackage{marginnote}
\usepackage{ulem}
\usepackage[yyyymmdd,hhmmss]{datetime}
\usepackage{ulem}
\usepackage{environ}
\usepackage{graphicx}
%\usepackage[top=Bcm, bottom=Hcm, outer=Ccm, inner=Acm, heightrounded, marginparwidth=Ecm, marginparsep=Dcm]{geometry}
\usepackage{courier}



\newtcolorbox{Example}[1]{colback=white,left=20pt,colframe=slideblue,fonttitle=\bfseries,title=#1}
\newtcolorbox{Solutions}[1]{colback=white,left=20pt,colframe=green,fonttitle=\bfseries,title=#1}
\newtcolorbox{Conseils}[1]{colback=white,left=20pt,colframe=slideblue,fonttitle=\bfseries,title=#1}
\newtcolorbox{Warning}[1]{colback=white,left=20pt,colframe=warning,fonttitle=\bfseries,title=#1}
\newtcolorbox{Note}[1]{colback=white,left=20pt,colframe=grey,fonttitle=\bfseries,title=#1}
\newtcolorbox{Dessin}[1]{colback=white,left=20pt,colframe=grey,fonttitle=\bfseries,title=#1}

\setlength\parindent{0pt}

  %Exercice environment
  \newcounter{exercice}
  \newenvironment{Exercice}[1][]
  {
  \par
  \stepcounter{exercice}\textbf{Question \arabic{exercice}:} (\faClock \enskip \textit{#1})
  }
  {\bigskip}
  

% Title
\newcommand{\titre}{\begin{center}
  \section*{Algorithmes et Pensée Computationnelle}
\end{center}}
\newcommand{\cours}[1]
{\begin{center} 
  \textit{#1}\\
\end{center}
  }



% \newcommand{\exemple}[1]{\newline~\textbf{Exemple :} #1}
\newcommand{\todo}[1]{\textcolor{red}{#1}}
%\newcommand{\attention}[1]{\newline\faExclamationTriangle~\textbf{Attention :} #1}

% Documentation url (escape \# in the TP document)
\newcommand{\documentation}[1]{\faBookOpen~Documentation : \href{#1}{#1}}

% Clef API
\newcommand{\apikey}[1]{\faKey~Clé API : \lstinline{#1}}
\newcommand{\apiendpoint}[1]{\faGlobe~Url de base de l'API \href{#1}{#1}}

%Listing Python style
\usepackage{color}
\definecolor{slideblue}{RGB}{33,131,189}
\definecolor{green}{RGB}{0,190,100}
\definecolor{blue}{RGB}{121,142,213}
\definecolor{grey}{RGB}{120,120,120}
\definecolor{warning}{RGB}{235,186,1}

\usepackage{listings}
\lstdefinelanguage{texte}{
    keywordstyle=\color{black},
    numbers=none,
    frame=none,
    literate=
           {é}{{\'e}}1
           {è}{{\`e}}1
           {ê}{{\^e}}1
           {à}{{\`a}}1
           {â}{{\^a}}1
           {ù}{{\`u}}1
           {ü}{{\"u}}1
           {î}{{\^i}}1
           {ï}{{\"i}}1
           {ë}{{\"e}}1
           {Ç}{{\,C}}1
           {ç}{{\,c}}1
           {ô}{{\^o}}1,
    columns=fullflexible,keepspaces,
	breaklines=true,
	breakatwhitespace=true,
}
\lstset{
  language=Python,
	%basicstyle=\bfseries\footnotesize,
  basicstyle=\footnotesize\ttfamily\small,
	breaklines=true,
	breakatwhitespace=true,
	commentstyle=\color{grey},
	stringstyle=\color{slideblue},
  keywordstyle=\color{slideblue},
	morekeywords={with, as, True, False, Fraction, Float, join, index, None, main, argparse, self, sort, __eq__, __add__, __mul__ ,__ne__, __radd__, __del__, __ge__, __gt__, split, os, endswith, is_file, scandir, @classmethod, this},
	deletekeywords={id},
	showspaces=false,
	showstringspaces=false,
	columns=fullflexible,keepspaces,
	literate=
           {é}{{\'e}}1
           {è}{{\`e}}1
           {ê}{{\^e}}1
           {à}{{\`a}}1
           {â}{{\^a}}1
           {ù}{{\`u}}1
           {ü}{{\"u}}1
           {î}{{\^i}}1
           {ï}{{\"i}}1
           {ë}{{\"e}}1
           {Ç}{{\,C}}1
           {ç}{{\,c}}1
           {ô}{{\^o}}1,
    numbers=left,
}

\newtcbox{\mybox}{nobeforeafter,colframe=white,colback=slideblue,boxrule=0.5pt,arc=1.5pt, boxsep=0pt,left=2pt,right=2pt,top=2pt,bottom=2pt,tcbox raise base}
\newcommand{\projet}{\mybox{\textcolor{white}{\small projet}}\xspace}
\newcommand{\optionnel}{\mybox{\textcolor{white}{\small Optionnel}}\xspace}
\newcommand{\advanced}{\mybox{\textcolor{white}{\small Pour aller plus loin}}\xspace}
\newcommand{\auto}{\mybox{\textcolor{white}{\small Auto-évaluation}}\xspace}


\usepackage{environ}
\newif\ifShowSolution
\NewEnviron{solution}{
  \ifShowSolution
	\begin{Solutions}{\faTerminal \enskip Solution}
		\BODY
	\end{Solutions}
  \fi}

\usepackage{environ}
\newif\ifShowExample
\NewEnviron{exemple}{
  \ifShowExample
	\begin{Example}{\faTerminal \enskip Exemple}
		\BODY
	\end{Example}
  \fi}

\usepackage{environ}
 \newif\ifShowConseil
  \NewEnviron{conseil}{
    \ifShowConseil
    \begin{Conseils}{\faLightbulb \quad Conseil}
      \BODY
    \end{Conseils}

    \fi}

    \usepackage{environ}
  \newif\ifShowWarning
  \NewEnviron{attention}{
    \ifShowWarning
    \begin{Warning}{\faTriangleExclamation \quad Attention}
      \BODY
    \end{Warning}

    \fi}

  \newif\ifShowNote
  \NewEnviron{note}[1]{
    \ifShowNote
    \begin{Note}{\faEye \quad Comment added by #1 at \date - \currenttime}
      \BODY
    \end{Note}

    \fi}  

    \newif\ifShowDessin
    \NewEnviron{dessin}[1]{
      \ifShowDessin
      \begin{Dessin}{\faPen \quad Figure}
        \BODY
      \end{Dessin}
  
      \fi}  

%\newcommand{\Conseil}[1]{\ifShowIndice\ \newline\faLightbulb[regular]~#1\fi}


\usepackage{listings}
\usepackage{array}
\newcolumntype{C}[1]{>{\centering\let\newline\\\arraybackslash\hspace{0pt}}m{#1}}
\lstset{upquote=true,basicstyle=\fontencoding{T1}\selectfont\ttfamily}
\hypersetup{
    colorlinks=true,
    linkcolor=blue,
    filecolor=magenta,      
    urlcolor=blue,
    pdfpagemode=FullScreen,
    }

\begin{document}

% Change the following values to true to show the solutions or/and the hints
\ShowSolutiontrue
\ShowConseiltrue
\ShowWarningtrue
\titre
\cours{Fonctions, mémoire et exceptions}

Le but de cette séance est d'approfondir vos connaissances en programmation. Au terme de cette séance, vous serez capable de: 
\begin{itemize}
    \item définir une fonction et l'utiliser dans un programme,
    \item utiliser des bibliothèques contenant des fonctions prédéfinies,
    \item connaître la portée d'une variable,
    \item comprendre comment fonctionne une pile d'exécution (call stack),
    \item gérer des exceptions.
\end{itemize}

\section{Opérations arithmétiques}
 \begin{Exercice}[3 minutes] [Java et Python] \textbf{Calculs (multiplication) }\\
 
    Créez deux variables nommées \texttt{facteur1} et \texttt{facteur2} en Java (\texttt{facteur\_1} et \texttt{facteur\_2} en Python), avec pour valeurs respectives 11 et 3. Multipliez la première variable par la deuxième et stockez le résultat dans une nouvelle variable \texttt{produit}. Vous pouvez afficher les différentes variables pour voir leurs valeurs. Vous pouvez répéter l'exercice avec l'addition et la soustraction. \\
    
     \begin{conseil}
           L'opérateur de multiplication est le *, celui de l'addition est le + et celui de soustraction est le -.
         
     \end{conseil}
     \begin{solution}
     
     \textbf{Python:}
     \lstinputlisting{solutions/question1.py}
     
     \textbf{\\Java:}
     \lstinputlisting{solutions/Question1.java}
            
     \end{solution}   
 \end{Exercice}
 
 \begin{Exercice}[10 minutes] [Java et Python] \textbf{Calculs (division)}\\
     Créez deux variables nommées \texttt{nbBonbons} et \texttt{nbPersonnes} en Java (\texttt{nb\_bonbons} et \texttt{nb\_personnes} en Python), avec pour valeurs respectives 11 et 3. Divisez la première variable par la deuxième et stockez le résultat dans une nouvelle variable nommée \texttt{bonbonsPersonnes} en Java (\texttt{bonbons\_personnes} en Python). Pour finir, calculez le nombre de bonbons restants via l'opérateur \% (modulo) et stockez le résultat dans une nouvelle variable \texttt{reste}. Vous pouvez afficher les différentes variables pour voir leurs valeurs. \\
     
      \begin{conseil}
         \begin{itemize}
            \item Attention: en Python il existe deux opérateurs de division, / effectue une division classique, tandis que // effectue une division avec arrondi inférieur.
            \item En Java, lorsque les deux opérandes sont de type int, l'opérateur / effectue une division entière. Si au moins l'un des deux opérandes est de type float ou double, il effectue une division en virgule flottante.
            \item L'opérateur \% (modulo) permet de calculer le reste d'une division. Par exemple, 10\%2=0.
         \end{itemize}
      \end{conseil}
      \begin{solution}
      
      \textbf{Python:}
      
      \lstinputlisting{solutions/question2.py}
      
      \textbf{\\Java:}
      \lstinputlisting{solutions/Question2.java}
             
      \end{solution}   
  \end{Exercice}
  
 \begin{Exercice}[5 minutes] [Java et Python] \textbf{Calculs (incrémentation / décrémentation)}\\
    Gardez vos variables de l'exercice précédent, augmentez de 1 la valeur de \lstinline|nbBonbons| en Java (\lstinline|nb_bonbons| en Python) et diminuez de 1 celle de \lstinline|nbPersonnes| en Java (\lstinline|nb_personnes| en Python).  \\
    
     \begin{conseil}
           Vous pouvez utiliser les opérateurs += et -= en Python, et les opérateurs ++ (incrémentation) et $--$ (décrémentation) en Java.
         
     \end{conseil}
     \begin{solution}
     
     \textbf{Python:}
     \lstinputlisting{solutions/question3.py}
     
     \textbf{\\Java:}
     \lstinputlisting{solutions/Question3.java}
            
     \end{solution}   
 \end{Exercice}

\section{Variables et Fonctions}
\begin{Exercice}[5 minutes] [Java \textbf{et} Python] \textbf{Les fonctions (fonctions basiques)}\\
    Définissez une fonction nommée \lstinline{ping()} qui, lorsqu'elle est appelée, affiche \texttt{pong}. \\
     
      \begin{conseil}
        Référez-vous aux diapositives du cours pour la création et l'appel des fonctions.    
      \end{conseil}
      \begin{solution}
      
          \textbf{Python}:
            \lstinputlisting{solutions/question4.py}
          
          \textbf{\\Java}:
            \lstinputlisting{solutions/Question4.java}
             
      \end{solution}   
  \end{Exercice}
  
    \begin{Exercice}[5 minutes] [Java et Python] \textbf{Les Fonctions (Fonction multiplication)}\\
    Définissez une fonction nommée \lstinline{multiplicateur()} qui prend deux arguments \textit{multiple1} et \textit{multiple2}, les multiplie et retourne le résultat. Stockez le résultat de \lstinline{multiplicateur(2,3)} dans une variable \textit{resultat} et affichez-la.   \\

    \begin{conseil}
        \begin{itemize}
            \item Référez-vous aux diapositives du cours pour la création et l'appel des fonctions.
            \item Pour retourner une valeur au lieu de l'afficher, utilisez le mot-clé \lstinline{return} (pour Python et Java).
        \end{itemize}        
    \end{conseil}
    \begin{solution}
        \textbf{Python}:
        \lstinputlisting{solutions/question5.py}
        
        \textbf{\\Java}:
        \lstinputlisting{solutions/Question5.java}
    \end{solution}   
\end{Exercice} 
  
  \begin{Exercice}[5 minutes] [Java et Python] \textbf{Les Fonctions (Fonctions Aire et Périmètre)}\\
      Définissez deux fonctions nommées \lstinline{aire()} et \lstinline{perimetre()} qui prennent un argument (\lstinline{rayon}) et renvoient respectivement l'aire et le périmètre d'un cercle. Stockez les résultats dans des variables \lstinline{aire_resultat} et \lstinline{perimetre_resultat} et affichez le contenu de ces variables.   \\
     
      \begin{conseil}
          \begin{itemize}
              \item Référez-vous aux diapositives du cours pour la création et l'appel des fonctions.
              \item Pour retourner une valeur au lieu de l'imprimer, utilisez le mot-clé \lstinline{return} (pour Python et Java).
              \item Pour rappel, le périmètre d'un cercle s'obtient en utilisant la formule $P = 2*\pi*r$ et l'aire s'obtient en utilisant la formule $A = r^2*\pi$.
          \end{itemize}        
      \end{conseil}
      \begin{solution}
          \textbf{Python}:
          \lstinputlisting{solutions/question6.py}
          
          \textbf{\\Java}:
          \lstinputlisting{solutions/Question6.java}
      \end{solution}   
  \end{Exercice} 

  \begin{Exercice}[5 minutes] [Python] \textbf{Portée des variables} \\
    Qu'affiche le programme suivant? \\

    \lstinputlisting{code/question7.py}

    
     \begin{conseil}
        \begin{itemize}
            \item Le mot-clé \lstinline{global} permet d'accéder aux variables globales (définies à l'extérieur de la fonction).
            \item Soyez attentif au type des éléments retournés par chacune des fonctions.
        \end{itemize}
     \end{conseil}
     \begin{solution}
        x = 2 est une variable globale (en dehors de la fonction), alors que x = 3 est une variable locale. \\
        La fonction imbriquée \lstinline|fonction2| fait appel à la variable globale, avec le mot-clé global devant.\\
        Lors du deuxième appel à \texttt{print}, il n'y a pas le mot clé global donc la variable locale est utilisée.\\
        La solution est donc :\\ 
        \lstinline{x=2}\\ 
        \lstinline{x=3}\\
        \lstinline[language={bash}]{None}\\
    \end{solution}
 \end{Exercice}

 \begin{Exercice}[10 minutes] [Java] \textbf{Portée des variables} \\
    Qu'affiche le programme suivant?
    \lstinputlisting{code/Question8.java}
    Que se passerait-il si on décommente la ligne 17 (c'est-à-dire si l'on enlève les //)?

    \begin{conseil}
        Tenir compte de l'endroit où chaque variable est définie ainsi que des opérations réalisées sur chaque variable.
    \end{conseil}

    \begin{solution}
        \begin{itemize}
            \item  
            \lstinline{a = 2} \\
            \lstinline{a = 2} \\
            \lstinline{b = 3} \\ 

            \item On obtiendra une erreur de compilation du programme car la variable \lstinline{b} est créée à l'intérieur de la méthode \lstinline{g()} et n'existe donc qu'au sein de celle-ci. La variable \lstinline{b} n'étant pas définie en dehors de la méthode \lstinline{g()}, celle-ci ne peut pas faire l'objet d'un appel en dehors de la méthode \lstinline{g()}.
        \end{itemize}
        \end{solution}
 \end{Exercice}


\section{Gestion des exceptions}

\begin{Exercice}[5 minutes] [Python] \textbf{Types d'erreurs} Qu'affiche le programme suivant?\\
    \lstinputlisting{code/question9.py}  
     \begin{conseil}
        Utilisée sur une chaîne de caractères, la fonction \lstinline{len()} renvoie le nombre de caractères de la chaîne. 
    \end{conseil}
     \begin{solution}
        \lstinline{Le résultat de la division est: 2.5}\\
        \lstinline{On ne peut pas diviser une chaîne de caractères.}
     \end{solution}   
 \end{Exercice}

\begin{Exercice}[5 minutes] [Python] \textbf{Types d'erreurs} Qu'affiche le programme suivant?\\
    \lstinputlisting{code/question10.py}
    
     \begin{conseil}
         La chaîne de caractères vide est représentée par ``'' 
    \end{conseil}
     \begin{solution}
        Nous ne pouvons pas diviser un nombre par 0.
     \end{solution}   
 \end{Exercice}

 \begin{Exercice}[5 minutes] [Python] \textbf{Types d'erreurs} Qu'affiche le programme suivant? \\
    \lstinputlisting{code/question11.py}
    
     \begin{conseil}
        \begin{itemize}
            \item \lstinline[language=python]|TypeError| survient lorsqu'une opération est effectuée sur des types incompatibles (exemple: \lstinline[language=python]|"abc" + 5|), tandis que \lstinline[language=python]|ValueError| survient lorsque le type est approprié, mais que la valeur est invalide (exemple: \lstinline[language=python]|int("abc")|).
            \item La fonction \lstinline{float()} permet de convertir une variable en nombre à virgule flottante.
            \item La fonction \lstinline{int()} permet de convertir une variable en nombre entier.
        \end{itemize}  
    \end{conseil}

     \begin{solution}
        \lstinline[language=python]|float(value2)| est valide car on transforme un nombre entier en un float. Aucune erreur ne va être lancée.\\
        \lstinline[language=python]|int(value1)| n'est pas valide si \lstinline[language=python]|value1| contient une chaîne qui ne représente pas un entier, comme \lstinline[language=python]|"abc"|. En revanche, une chaîne comme \lstinline[language=python]|"42"| peut être convertie en entier. \\
        
        La chaîne \texttt{"Algorithms"} provoque ainsi une erreur avec le message suivant:\\
        Nous ne pouvons pas convertir une chaîne de caractères en nombre
     \end{solution}

\end{Exercice}
 
\begin{Exercice}[5 minutes] [Java] \textbf{Gestion d'erreurs} 
    Le programme suivant est incorrect. Que devez-vous modifier pour qu'il fonctionne correctement?
    \\
    \begin{conseil}
        Référez-vous à la diapositive 31 du cours de cette semaine.           
     \end{conseil}
    \lstinputlisting{code/Question12.java}
     
     \begin{solution}
        À l'intérieur de la méthode \lstinline{division}, lorsque \lstinline{b} est égal à 0, le programme lève une exception de type \lstinline{ArithmeticException}. Mais dans la méthode \lstinline{main}, le bloc catch est configuré pour intercepter l'exception \lstinline{IndexOutOfBoundsException}. Pour corriger cette erreur, une des solutions est de remplacer \lstinline{IndexOutOfBoundsException} par \lstinline{ArithmeticException}.
     \end{solution}   
 \end{Exercice}

\begin{Exercice}[5 minutes] [Java] \textbf{Gestion d'erreurs} 
    \begin{enumerate}
    \item Qu'affiche le programme suivant, en supposant que l'utilisateur saisit les entrées suivantes? \\

    \begin{itemize}
        \item Entrez le premier entier: 12
        \item Entrez le deuxième entier: 2
        \item Entrez l'opération (+, -, *, /): *
    \end{itemize}

    \begin{conseil}
            Consultez \href{https://moodle.unil.ch/mod/resource/view.php?id=2218920f}{ce document} sur le \texttt{switch} en Java. Référez-vous à la diapositive 23 du cours. 
     \end{conseil}

     \item Que devez-vous modifier pour qu'il fonctionne?
    \end{enumerate}
    \lstinputlisting{code/Calculator.java}
    \begin{solution}
    \texttt{Entrez le premier entier: 12 \\ 
Entrez le deuxième entier: 2 \\
Entrez l'opération (+, -, *, /): * \\ 
24.0\\
6.0\\
Exception in thread "main" java.lang.IllegalArgumentException: L'opération n'est pas valide \\
        at Calculator.main(Calculator.java:28)}
    \end{solution}
\begin{solution}
Il faut simplement ajouter un \texttt{break;} après l'affichage du résultat d'une opération. 
\end{solution}

 \end{Exercice}

 \begin{Exercice}[5 minutes]
   Qu'affiche le programme suivant, en supposant que l'utilisateur exécute le programme de la manière suivante dans le terminal? \\ \newline 
   \texttt{python3 question14.py 9} \\ 
   \begin{conseil}
   Référez-vous à la diapositive 15 du cours.
   \end{conseil}
   \lstinputlisting{code/question14.py}
   \begin{solution}
   
\texttt{Une erreur est survenue \\
Valeur finale de balance: 20
    } \\ \\
    Cette erreur est due à la soustraction d'une chaîne de caractères (\texttt{sys.argv[1]}) d'un nombre, ce qui n'est pas définie. 
   \end{solution}
 \end{Exercice}

 \begin{Exercice}[5 minutes]
Laquelle des affirmations suivantes est \textbf{fausse} concernant la différence entre une Error (erreur) et une Exception (exception) en Java?
\begin{conseil}
    Référez-vous à la diapositive 36 du cours.
\end{conseil}
\begin{enumerate}[label=(\Alph*)]
    \item Les erreurs sont généralement causées par l'environnement dans lequel l'application s'exécute, tandis que les exceptions sont généralement causées par l'application elle-même.
    \item Les erreurs sont non vérifiées (\textit{unchecked}) et n'ont pas besoin d'être déclarées dans la clause \texttt{throws} d'une méthode, tandis que les exceptions peuvent être soit vérifiées (\textit{checked}), soit non vérifiées (\textit{unchecked}).
    \item Les erreurs devraient généralement être interceptées et gérées par l'application afin d'assurer une récupération fluide, tandis que les exceptions peuvent être récupérables ou non.
    \item Les erreurs indiquent généralement des problèmes graves qu'une application raisonnable ne devrait pas essayer d'intercepter.
    \item \texttt{RuntimeException} sont des exceptions non vérifiées (\textit{unchecked}).
\end{enumerate}

\begin{solution}
(C) L'affirmation C est fausse car elle inverse l'approche habituelle. 
Les erreurs (telles que \texttt{OutOfMemoryError}, \texttt{StackOverflowError}) ne doivent généralement pas 
être détectées et gérées par les applications, car elles indiquent des problèmes graves qui ne peuvent 
généralement pas être résolus. Les exceptions, en revanche, sont conçues pour être détectées et gérées par l'application lorsque cela est approprié, 
car elles représentent souvent des conditions récupérables (telles que \texttt{ArithmeticException} ou \texttt{IllegalArgumentException}).
\end{solution}
 \end{Exercice}

\end{document}
